\documentclass[12pt,xcolor={dvipsnames,x11names},t]{beamer}
\usepackage{engin310}

\title{ENGIN310: Programación y Algoritmos}
\subtitle{Semestre de otoño de 2017}
\author{Nicolás Eugenio Ozimica\\[.4em]\texttt{<nozimica@fen.uchile.cl>}}

\date{Apuntes de clase}

\usepackage[export]{adjustbox}

\begin{document}

%----------- titlepage ----------------------------------------------%
\begin{frame}[plain]
    \centering \includegraphics[height=26mm]{img/logoFen.png}\\
    \centering
    \titlepage
\end{frame}

% \providecommand{\tightlist}{%
%     \setlength{\itemsep}{0pt}\setlength{\parskip}{0pt}}
\def\opacityblock{.3}

\tikzstyle{simplestep} = [normalfig, rectangle, text width=7em, bottom color=orange!30, draw=orange!90!black!50, rounded corners=2mm]
\tikzstyle{na} = [baseline=-0.5ex,remember picture]
\tikzstyle{arrowInList} = [draw=blue,thick,->]

\def\tightlist{}

\section{Introducción al curso}
\subsection{Algoritmos}
\include{tex/1.01algoritmos}
\subsection{Variables y tipos de dato}
\include{tex/1.02variablesytipos}
\section{Introducción a Python}
\subsection{Expresiones simples}
\include{tex/2.01simpleexpr}
\subsection{Variables en Python}
\include{tex/2.11simplevars}
\subsection{Funciones}
\include{tex/2.21tareasrepetidas}
\include{tex/2.24funciones}
\subsection{Receta de diseño}
\include{tex/2.27recetadisenno}
\end{document}
